\setcounter{aufgabe}{28}
\einheitenkopf{Einheit 4 von 5 · Kenngrößen}

\begin{aufgabe}{Kenngrößen -- welche Kenngröße? Kreuze an, wonach gefragt wird. Rechne nicht.}
\begin{teile}
\teil „Wie weit liegen der kleinste und der größte Wert auseinander?“ \\ \kreuz{Spannweite}\kreuz{Modalwert}\kreuz{Median}\kreuz{Mittelwert}
\teil „Welcher Wert kommt am häufigsten vor?“ \\ \kreuz{Spannweite}\kreuz{Modalwert}\kreuz{Median}\kreuz{Mittelwert}
\teil „Welcher Wert steht in der Mitte, wenn man alle Werte ordnet?“ \\ \kreuz{Spannweite}\kreuz{Modalwert}\kreuz{Median}\kreuz{Mittelwert}
\teil „Wie viele Besucher kamen im Durchschnitt?“ \\ \kreuz{Spannweite}\kreuz{Modalwert}\kreuz{Median}\kreuz{Mittelwert}
\teil „Welche Schuhgröße wurde am öftesten verkauft?“ \\ \kreuz{Spannweite}\kreuz{Modalwert}\kreuz{Median}\kreuz{Mittelwert}
\end{teile}
\end{aufgabe}

\begin{aufgabe}{Kenngrößen -- ordnen. Schreibe die Werte der Größe nach geordnet auf. Bestimme noch nichts.}
\begin{teile}
\teil $7 \quad 3 \quad 9 \quad 4 \quad 6$ \\[10pt]
\teil $12 \quad 9 \quad 15 \quad 9 \quad 11 \quad 14$ \\[10pt]
\teil $2{,}5 \quad 1{,}8 \quad 3{,}1 \quad 2{,}0$ \\[10pt]
\teil $105 \quad 98 \quad 112 \quad 98 \quad 103$ \\[10pt]
\teil $0{,}7 \quad 1{,}2 \quad 0{,}9 \quad 1{,}5 \quad 0{,}8$ \\[10pt]
\end{teile}
\end{aufgabe}

\begin{aufgabe}{Kenngrößen -- Spannweite, Modalwert, Median. Gib bei a) bis d) nur Minimum und Maximum an.}
\beispiel{\text{Werte } 5,\ 9,\ 3,\ 9,\ 7: \quad &\text{Minimum } 3, \quad \text{Maximum } 9}
\begin{teile}
\teil $4 \quad 11 \quad 7 \quad 2 \quad 9$ \hfill \feld{\text{Min}} \feld{\text{Max}}
\teil $21 \quad 15 \quad 30 \quad 18$ \hfill \feld{\text{Min}} \feld{\text{Max}}
\teil $3{,}4 \quad 2{,}9 \quad 4{,}1 \quad 3{,}0$ \hfill \feld{\text{Min}} \feld{\text{Max}}
\teil Die Tagestemperaturen einer Woche: $9$, $12$, $7$, $14$, $11$, $12$, $8$ (in $^\circ$C) \hfill \feld{\text{Min}} \feld{\text{Max}}
\teil Wie groß ist die Spannweite der Werte aus a)? \leerfeld
\teil Wie groß ist die Spannweite von $8{,}7 \quad 9{,}2 \quad 8{,}5 \quad 9{,}0$? \leerfeld
\teil Wie lautet der Modalwert von $6 \quad 4 \quad 6 \quad 9 \quad 4 \quad 6$? \leerfeld
\teil Wie lautet der Median von $5 \quad 2 \quad 8 \quad 1 \quad 9$? \leerfeld
\teil Wie lautet der Median von $14 \quad 11 \quad 20 \quad 17$? \leerfeld
\teil Lies im Diagramm den kleinsten und den größten Niederschlagswert ab und gib die Spannweite an.
\par \feld{\text{Min}} \feld{\text{Max}} \feld{\text{Spannweite}}
\end{teile}

\noindent\saeulen{Jan/45,Feb/30,Mär/60,Apr/20,Mai/55}{70}{10}{Niederschlag in mm}
\end{aufgabe}

\begin{aufgabe}{Kenngrößen -- Mittelwert. Berechne das arithmetische Mittel. Ab e) darfst du den Taschenrechner benutzen.}
\beispiel{4,\ 6,\ 8,\ 10: \quad (4+6+8+10) : 4 &= 28 : 4 = 7}
\begin{teilezwei}
\tz $3 \quad 5 \quad 7$ \leerfeld & \tz $2 \quad 4 \quad 6 \quad 8$ \leerfeld \\
\tz $10 \quad 20 \quad 30$ \leerfeld & \tz $1 \quad 2 \quad 3 \quad 4 \quad 5$ \leerfeld \\
\end{teilezwei}
\begin{teile}
\teil $12{,}5 \quad 13{,}5 \quad 14{,}0 \quad 12{,}0$ \leerfeld
\teil Sechs Messwerte: $8{,}4 \quad 9{,}1 \quad 7{,}6 \quad 8{,}9 \quad 9{,}0 \quad 8{,}4$ -- runde auf eine Stelle nach dem Komma. \leerfeld
\teil Ein Kino zählte in 17 Vorstellungen zusammen 4390 Besucher. Wie viele Besucher kamen im Durchschnitt je Vorstellung? Runde sinnvoll. \leerfeld
\teil Acht Werte haben zusammen die Summe 344. Wie groß ist ihr Mittelwert? \leerfeld
\teil Bei sechs Messungen beträgt der Mittelwert 15. Fünf der Werte sind $12$, $18$, $14$, $11$ und $16$. Wie groß ist der sechste Wert? \leerfeld
\teil In einer Gruppe sind neun Personen. Ihr Alter in Jahren: $16$, $28$, $31$, $35$, $38$, $41$, $30$, $35$, $52$. Berechne das Durchschnittsalter. Die älteste und die jüngste Person verlassen die Gruppe. Berechne das Durchschnittsalter der übrigen sieben Personen und begründe, warum es sich nicht ändert. (P10 2025)
\par \feld{\text{neun Personen}} \quad \feld{\text{sieben Personen}} \\[12pt]
\end{teile}
\end{aufgabe}

\begin{aufgabe}{Kenngrößen -- Aussagen prüfen. Zu den Werten $9 \quad 4 \quad 7 \quad 4 \quad 12$ werden vier Aussagen gemacht. Kreuze an, ob sie stimmen, und schreibe hinter die falschen Aussagen den richtigen Wert.}
\begin{teile}
\teil „Die Spannweite ist 8.“ \hfill \janein \quad richtig: \leerfeld
\teil „Der Modalwert ist 12.“ \hfill \janein \quad richtig: \leerfeld
\teil „Der Median ist 7.“ \hfill \janein \quad richtig: \leerfeld
\teil „Der Mittelwert ist 7.“ \hfill \janein \quad richtig: \leerfeld
\end{teile}
\end{aufgabe}

\begin{aufgabe}{Kenngrößen -- Fehler finden.}
Zu den Werten $12 \quad 5 \quad 9 \quad 3 \quad 7$ gibt Nina an:
\rechnung{\text{Median} &= 9}
\begin{teile}
\teil Finde den Fehler, benenne ihn und gib den richtigen Median an. \\[12pt]
\teil Bestimme den Median von $20 \quad 14 \quad 25 \quad 11 \quad 18 \quad 22$. \leerfeld
\end{teile}
\end{aufgabe}

\begin{aufgabe}{Kenngrößen -- Begründen. Rechne nicht.}
\begin{teile}
\teil In einem Betrieb verdienen neun Angestellte je 2400\,€ im Monat, die Chefin 12\,000\,€. Erkläre, warum der Median den Verdienst einer „normalen“ Angestellten besser beschreibt als der Mittelwert. \\[12pt]
\teil Zu acht Werten wird ein neunter Wert hinzugefügt, der genau so groß ist wie der bisherige Mittelwert. Ändert sich der Mittelwert? Begründe. \\[12pt]
\end{teile}
\end{aufgabe}
